\chapter{Cài đặt}

Pulse chạy trên cụm Kubernetes 6 node của phòng lab. Mỗi worker có một loader eBPF viết bằng C và một worker Python chấm điểm. Bridge cảnh báo và RCA chạy trên master đầu tiên. Code, protocol và biên nhận nằm trong cùng một repo git.

\section{Cụm thử nghiệm}

\begin{longtable}[]{@{}>{\raggedright\arraybackslash}p{\dimexpr 0.271\linewidth-2\tabcolsep\relax}>{\raggedright\arraybackslash}p{\dimexpr 0.729\linewidth-2\tabcolsep\relax}@{}}
\toprule\noalign{}
Thành phần & Cấu hình \\
\midrule\noalign{}
\endhead
\bottomrule\noalign{}
\endlastfoot
Node & 3 master (control plane), 3 worker; VM trên ESXi \\
Phần mềm & Kubernetes v1.34.10, Ubuntu 24.04, kernel 6.8.0, Cilium, Istio ambient \\
Ứng dụng & AIMS (An Internet Media Store): frontend, API gateway, các service Python, PostgreSQL (CloudNativePG), Kafka, RabbitMQ, Redis, MinIO; 21 workload, khoảng 80 pod \\
Tải normal & Bốn generator trong namespace \texttt{production}, chuyển qua các chế độ steady, toolmix, peak, burst, recovery theo lịch đăng ký \\
Công cụ runtime có sẵn & Tetragon (nguồn sự kiện cho RCA), Falco (baseline luật) \\
Lưu trữ & NFS riêng cho dữ liệu các lượt chạy \\
\end{longtable}

\section{Thu thập trong kernel}

\begin{itemize}
\tightlist
\item
  \textbf{Loader eBPF} (\texttt{pulse\_\allowbreak{}counter\_\allowbreak{}extended}) gắn chương trình vào \texttt{raw\_\allowbreak{}tp/\allowbreak{}sys\_\allowbreak{}enter} và \texttt{fexit/\allowbreak{}\_\allowbreak{}\_\allowbreak{}seccomp\_\allowbreak{}filter}. Nhờ vậy nó đếm cả syscall bị seccomp từ chối.
\item
  \textbf{Map theo cgroup} giữ ba loại bộ đếm: số đếm theo từng ID syscall (0--1023), 64 bin ID và 64 bin cặp syscall. Bin cặp dùng syscall trước đó của từng thread, nên kernel phải nhớ syscall cuối của mỗi thread.
\item
  \textbf{Danh sách cgroup cho phép} do dịch vụ resolver ghi. Resolver ánh xạ Pod UID → container → cgroup → phiên bản workload.
\item
  \textbf{Cứ 500 ms,} loader đọc map và in một frame JSON ra stdout. Mỗi frame có bộ đếm tích lũy của mọi cgroup được theo dõi.
\item
  \textbf{Khóa bằng proof:} loader và object được khóa bằng SHA-256 cùng một “proof” gắn với phiên bản kernel. Worker từ chối chạy trên kernel khác.
\end{itemize}

\section{Worker chấm điểm}

Worker (\texttt{pipeline/\allowbreak{}matched\_\allowbreak{}worker.\allowbreak{}py}) chạy như một unit systemd trên mỗi worker, giới hạn 3 core và 3 GB RAM. Mỗi frame đi qua các bước:

\begin{enumerate}
\def\labelenumi{\arabic{enumi}.}
\tightlist
\item
  Kiểm tính toàn vẹn: tổng theo ID phải bằng tổng của probe, và bin băm phải bằng phép chiếu của số đếm theo ID. Frame sai bị ghi vào \texttt{issues} và lịch sử được xóa.
\item
  Tính hiệu số với frame trước, dựng 249 feature và ngữ cảnh 4 cửa sổ.
\item
  Chấm theo batch: gom mọi ngữ cảnh cùng workload trong một frame, tối đa 256 bản ghi.
\item
  Áp chính sách, chiết khấu theo replica và tripwire, rồi ghi decision.
\end{enumerate}

\textbf{Đường chấm nhanh} (\texttt{detector/\allowbreak{}fast\_\allowbreak{}forest.\allowbreak{}py}) làm đúng phép tính của scikit-learn: cộng xác suất của từng cây theo thứ tự, bắt đầu từ số 0 float64, nhưng bỏ phần chi phí Python của mỗi cây.

\begin{itemize}
\tightlist
\item
  Batch đầu tiên của mỗi model được chấm bằng cả hai đường. Chỉ cần lệch một giá trị là model đó quay về đường cũ suốt tiến trình.
\item
  Replay 41.250 dòng feature live qua đường tham chiếu cho 0 dòng lệch.
\item
  CPU của worker giảm từ khoảng 94\% xuống 13--31\% một core.
\end{itemize}

\textbf{Khóa bản đóng băng.} Trước khi chấm, worker so SHA-256 của 13 file code chấm điểm với một snapshot đóng băng, và kiểm hash của model bundle cùng chính sách. Nhờ vậy mọi kết quả live gắn với đúng một phiên bản code.

\section{Trao đổi giữa worker}

Mỗi worker mở một cổng TCP (9479), phát mặt nạ ID syscall của từng cửa sổ và nhận mặt nạ của hai worker còn lại. Tin của peer tới sau khoảng 20 ms, nên bucket mới nhất có thể chưa đủ lúc quyết định.

Cùng kênh này mang cửa sổ tham chiếu cho tripwire. Trong lượt cuối, không cửa sổ tham chiếu nào của peer đến muộn.

\section{Ghi dữ liệu ngoài vòng chấm điểm}

Các lượt đầu ghi feature và decision thẳng lên NFS trong vòng chấm điểm. Mỗi lần NFS chậm là một lần frame bị trễ; có frame trễ tới 31,7 s. Hệ thống cuối sửa theo ba cách:

\begin{itemize}
\tightlist
\item
  \textbf{Ghi nền:} mỗi file chia sẻ được ghi qua một hàng đợi và thread riêng. Thread này \texttt{fdatasync} và bỏ page cache mỗi 30 s.
\item
  \textbf{Bản decision trên đĩa local} của worker cho bridge, RCA và campaign đọc.
\item
  \textbf{Trạng thái ghi không đồng bộ.}
\end{itemize}

Kết quả ở lượt cuối:

\begin{itemize}
\tightlist
\item
  độ trễ frame lớn nhất còn 2,4 s;
\item
  lần ghi nền lâu nhất 35 ms, không lần nào phải chờ hàng đợi;
\item
  số lần cách nhịp còn 4, 2 và 1 trên ba worker, so với 8--10 ở lượt trước.
\end{itemize}

\section{Bridge và RCA}

\begin{itemize}
\tightlist
\item
  \textbf{Bridge cảnh báo} chạy trên master, bám file decision local của ba worker và đẩy đợt cảnh báo vào Alertmanager. Ở lượt cuối có 3 push lỗi, cả 3 được thử lại thành công; 5.840/5.840 đợt tới nơi.
\item
  \textbf{Subscriber RCA} đọc luồng sự kiện Tetragon của ba worker vào ring SQLite local, dựng case theo nhịp 0,25 s và giữ tối đa 20.000 case.

  \begin{itemize}
  \tightlist
  \item
    Luồng \texttt{tetra\ getevents} được xoay vòng có giới hạn thời gian sống.
  \item
    DaemonSet Tetragon chạy với \texttt{shareProcessNamespace} để các process \texttt{timeout} kết thúc được thu dọn. Trước khi sửa có 3.931 process zombie trên ba worker.
  \end{itemize}
\end{itemize}

\section{Calibration tự làm mới trong worker}\label{sec:live-recal}

Calibration của bundle cũ dần theo ngày (mục~\ref{sec:limits-calibration}). Worker vì vậy có thêm một chế độ làm mới theo chu kỳ, \texttt{rolling-\allowbreak{}12-\allowbreak{}1}: mỗi giờ, hai mảng calibration của từng workload được thay bằng điểm của các cửa sổ eligible trong 12 giờ trước, gộp cả ba worker. Đây đúng là chiến lược đã cho kết quả tốt nhất trên soak 48 giờ ở phân tích offline.

\begin{itemize}
\tightlist
\item
  \textbf{Không đổi file chấm điểm:} cả đường chấm thường lẫn đường chấm nhanh đọc hai mảng \texttt{calibration\_scores} và \texttt{support\_calibration} của model ở mỗi lần gọi. Vì vậy chỉ cần thay hai mảng giữa hai frame; estimator, ranh giới support, chính sách và 13 file bị ghim SHA-256 giữ nguyên.
\item
  \textbf{Dữ liệu đi ngoài vòng chấm điểm:} mỗi worker ghi điểm thô của mọi cửa sổ vào một bộ đệm theo giờ trên đĩa local (31 byte mỗi cửa sổ). Một thread nền chép giờ vừa xong lên NFS sau 60 s, rồi sau 300 s đọc 12 giờ của cả ba worker và dựng mảng mới. Vòng chấm điểm chỉ nhận mảng đã dựng và thay vào giữa hai frame.
\item
  \textbf{Ghi vết:} mỗi decision mang số \texttt{calibration\_epoch}; epoch chỉ tăng khi có ít nhất một workload được thay. Workload chưa đủ 1.200 cửa sổ giữ mảng cũ.
\item
  \textbf{Khoảng campaign:} trong lượt đánh giá, các khoảng của campaign được loại khỏi calibration, đúng như bản offline. Hệ thống vận hành thật không biết lúc nào bị tấn công, nên đây là một giới hạn (mục~\ref{sec:limits-calibration}).
\item
  \textbf{Kiểm cài đặt:} một script dựng lại từng calibration đã áp dụng chỉ bằng các thành phần offline, rồi so mọi p-value live dưới epoch đó. Trong lượt smoke 75 phút (29/10), cả ba worker áp dụng epoch 1 cho 21/21 workload; 62.887 dòng decision khớp bản dựng lại 100\%, và kích thước mảng giống nhau trên ba worker.
\end{itemize}

Lượt xác nhận 24 giờ theo protocol \texttt{live-\allowbreak{}recalibration-\allowbreak{}confirmation-\allowbreak{}v1} so nhánh live R với chính các dòng đó tính lại bằng calibration của bundle (F). Luật thắng: gánh thấp hơn ít nhất 10\%, recall không kém quá 2 điểm, control không cao quá 1 điểm và p95 không chậm quá 0,2 s.

\section{Quy trình một lượt live}

\begin{enumerate}
\def\labelenumi{\arabic{enumi}.}
\tightlist
\item
  \textbf{Planner} tạo một file \texttt{CONFIG.\allowbreak{}json} bất biến. File này chứa hash của model, chính sách và fixture, lịch tải, lịch trial đã xáo trộn bằng seed cố định, và cấu hình từng worker.
\item
  \textbf{Launcher} cài worker lên ba node, khởi động RCA và bridge, rồi đến traffic và campaign.
\item
  \textbf{90 phút đầu là pha chọn:} tham chiếu của tripwire và baseline binary của RCA được dựng ra rồi đóng băng. Tấn công chỉ bắt đầu sau mốc này 120 s.
\item
  \textbf{Sau khi capture kết thúc,} một chuỗi phân tích tự chạy theo protocol đã commit. Biên nhận JSON được chép vào repo.
\end{enumerate}

Không job nặng nào chạy cùng lúc với lượt live, để điều kiện đo không đổi giữa các lượt.
